\section{Conclusions}\label{sec:conclusion}
Limited retained information creates a service-design frontier, and participation timing changes that frontier. With approval before a private draw, the provider can improve a restricted codebook by using adjacent lotteries and optimizing its highest level continuously. With realization-by-realization caps, randomization cannot improve any budget in the stated renewal architecture. A bounded-overrun agreement connects these alternatives in an exactly solved example, while explicit bit prices show when the institution changes the optimal writable capacity.

The model-specific optimization step is the cap-anchored reduction and the preservation of Monge structure after continuous terminal minimization. Its favorable one-sided domains then admit classical matrix search, yielding complete exact frontiers with linear rational work per budget layer. The resulting values support a stated implementation menu that prices writable bits, read-only representation, evaluation work, and compilation separately. Those distinctions are essential: neither a small symbol alphabet nor a better asymptotic search bound alone establishes a cheaper deployed service system.

\paragraph{Code, data, and reproducibility.}
The accompanying repository contains one current code directory and one reproduction entry point, \texttt{code/reproduce.py}, within the R37 revision package. It supplies exact engine snapshots, the attributed PADS excerpt, direct-network HiGHS comparisons, tests, and machine-readable measurements. The root article and electronic companion are the current readers. The preservation map locates every inherited result and the exact preceding readers; old theorem modules and evidence are unchanged. A build manifest binds the current sources, measurements, reader hashes, page counts, and whole-tree preservation checks to a scientific source commit. All experiment inputs are synthetic and public. No journal submission or editorial acceptance is represented by the repository build.
